% Modernised reading — fonds Alexandre Grothendieck, shelfmark 130
% (pages 1-49, the whole folder).
%
% This is an INTERPRETATION, not a transcription. It re-reads the pages in
% today's notation and vocabulary, states the hypotheses the manuscript leaves
% implicit, and drops the critical apparatus so that the mathematics reads
% straight through. Everything the brackets and underlines carried moves into a
% footnote. It is held to being mathematically correct as it stands; where the
% manuscript is loose or mistaken, the true statement is what appears here and
% the footnote says what the page has. In any dispute of substance the
% transcription governs: whoever cites Grothendieck must cite batch-01.fr.tex
% (pp. 1-20), batch-02.fr.tex (pp. 21-40) or batch-03.fr.tex (pp. 41-49).
%
% Pass: Opus 5.5 (claude-opus-5-5), 2026-10-03 — first pass, unchecked against the pages by a human.
% The folder was read whole, its three batches together; this is one document
% for the shelfmark. It was made on Opus 5.5 and has not been measured against
% the reference reading of folder 115 (made on Opus 5). The literature named
% in the footnotes is cited from memory and was not checked against the
% sources.
%
% Scope. Page 1 is a cover (« Formules de Néron »); pages 2, 8, 15, 20 and 33
% are blank. Pages 7, 17 and 19 belong to a Bourbaki typescript (n° 383, on
% root systems); of them only the five pencil lines in the margin of page 7
% are taken up, and the transcription does not assure they are his. Pages 22,
% 24, 26, 28, 30 and 32 are typed pages of an unrelated duplicated typescript;
% page 49 a typed page of someone else's text (n° 357), crossed out. None of
% these is described here beyond identification. The dating « [vers 1971] » is
% the inventory's for the folder; the inventory group « Descentes » in which
% the folder sits is dated « [avant 1970] » (observation of batch-01's header,
% recorded, not resolved).
%
% Spine. Four runs, linked by one theme (local symbols and heights), kept
% distinct.
%   pp. 1-7    « Formules de Néron »: Néron's local symbol (X, a)_v on an
%              abelian variety over the fraction field of a DVR, split as a
%              correction v(X, a) in (1/m)Z plus intersections on the Néron
%              model; pairing A*(K) x A(K) -> Z/mZ; global sum = Néron-Tate
%              (up to sign); p. 6: closure of Pic^0 in Pic_{X/S} as a Néron
%              model.
%   pp. 9-14   a surface X, a pencil of hyperplane sections, the blow-up X~ ->
%              P^1: kernel and image of NS -> J'(k(t)), the formula
%              rho(X) + z - 2 = rank J'(k(t)) (Shioda-Tate shape), and the
%              question Q(i xi, i eta) = - xi.eta (Shioda's height formula).
%   pp. 16-31  the local symbol <C_eta, D_eta>_v of cycles alg. equivalent to
%              0 should be i(C.D) on a regular model; properties; reduction,
%              when D is a divisor and C a 0-cycle, through the Albanese map
%              to the abelian case of pp. 3-5 and the normalisation by div f.
%              Breaks off on the non-smooth case (p. 31).
%   pp. 34-48  « Quasi-fonctions de Néron-Weil »: M_K, M_K-bounded sets,
%              admissible functions, the sheaf Phi_X, the sheaf of
%              quasi-functions QF_X, 0 -> Phi -> QF -> Div -> 0, kernel U of
%              M_K-units, global sections, surjectivity onto Pic for X
%              quasi-projective reduced to O(1) on P^r. Pages 45 and 47 are
%              cancelled and form their own thread; p. 44 continues on 46 and
%              46 on 48.
%
% Notation fixed for the folder: v is an additive (logarithmic) valuation;
% \mathcal{A} the Néron model (his A, then underlined A on p. 23 and 31);
% A^\vee for his A^* (dual abelian variety); a, b for his underlined 0-cycles;
% D_a, D_l (divisors alg., lin. equivalent to 0), Z_0^0, Z_l (0-cycles of degree
% 0, of sum 0); \langle C, D \rangle_v for the local symbol of pp. 16-31
% (his ( , )_v on p. 18); \Phi_X for his hatted / underlined Phi; QF_X;
% U the M_K-units. Sheaf underlines dropped.
%
% Departures from the pages (each footnoted in the body):
%  p. 3    0-cycles taken of degree 0, Z_l = kernel of the sum map (ours);
%          residue-field finiteness not needed with v normalised.
%  p. 5    the global pairing is A^\vee(K) x A(K) (page: A*(K) x A*(K));
%          values in (1/m)Z need integral weights and m = lcm of the m_v, and
%          fail for number fields (archimedean places); « forme de Néron-Tate »
%          holds up to sign (p. 21's sign is the right one with v additive).
%  p. 6    separatedness of the closure and an étale-local section supplied
%          (his margin points at the second); statement in dim S = 1.
%  p. 11   « extension pure de k(eta) » read « de k »; Pic^0(X) -> J^fixe(k)
%          asserted isomorphic, kept only up to isogeny (enough for ranks).
%  p. 13   z must be the number of base points card Z = (Y_0.Y_0), not of
%          singular sections (uncertain reading would make the formula false).
%  p. 14   the page asks; the answer (yes, with all fibres irreducible, up to
%          the normalisation of Q) is ours, cited from Shioda.
%  p. 16   cycles indexed by codimension (page: « dimensions », then
%          « codimensions »; the sums agree); model-independence is the
%          programme, not a result; needs a vertical correction in general.
%  p. 29   « sans composante horizontale » corrected to « verticale ».
%  p. 31   « point singulier de X' » read « de X_0 ».
%  p. 42   theta(f) = (-phi_f, f): the exponent of f is illegible; ours.
%  p. 46   « remplacer H^1 par 0 » kept as « the maps to H^1 vanish ».
%  p. 48   the P^r case is completed in a footnote (ours).
%
% Announced and never established: the « Lemme » of p. 29 (illegible); the
% Cor. of p. 5 (stated, not proved); the « à examiner » of p. 4; the margin
% « se généralise aussitôt aux variétés quelconques » (p. 4) beyond pp. 25-31;
% the formula of p. 14; model-independence of <C, D>_v (p. 16); the
% properties (iv) continuity and (v) base change (p. 18), only named; the
% non-smooth case b) (p. 31); the point « à vérifier » of p. 41 (admitted);
% whether the condition of p. 45 suffices; the P^r case of p. 48; the
% question of p. 37 on bounded heights.
\documentclass[11pt,a4paper]{article}
\input{../preamble/grothendieck}

\folder{130}
\pages{1}{49}
\dating{[vers 1971]}
\watermark{Édition de démonstration\\interprétation personnelle de l'œuvre}
\foldertitle{Hauteurs et symboles locaux : notes manuscrites (s.d.) — lecture modernisée du dossier entier}

\begin{document}

\section*{Résumé}

\begin{resume}
Ces feuillets cherchent à mesurer la « distance » entre deux objets
géométriques définis sur un corps de nombres, ou sur un corps de fonctions,
et à la décomposer en une somme de contributions locales, une par nombre
premier.

Le modèle est élémentaire. La taille d'un nombre rationnel $a/b$ en
écriture réduite se lit nombre premier par nombre premier : on regarde
combien de fois chaque $p$ divise $a$ et $b$, et l'on additionne ces
informations locales. Sur une courbe ou une variété, les points rationnels
ont de même une « hauteur », qui mesure leur complexité arithmétique, et
André Néron avait montré qu'elle se découpe en « symboles locaux », un par
place du corps. Le problème que se posent ces pages est de comprendre ce
que sont ces symboles locaux en géométrie : pour une place donnée, on
regarde la variété « réduite modulo $p$ », c'est-à-dire un modèle sur
l'anneau local, et l'on espère que le symbole local n'est autre qu'un
nombre d'intersection sur ce modèle — le nombre de fois que deux objets se
rencontrent quand on les réduit modulo $p$.

L'idée centrale est que cela est presque vrai. Sur une variété abélienne,
Grothendieck sépare le symbole de Néron en deux parts : un nombre
d'intersection sur le modèle de Néron, entier, et un terme correctif qui
ne dépend que de la façon dont la réduction modulo $p$ se casse en
plusieurs morceaux ; ce terme est une fraction de dénominateur le nombre
de morceaux. D'où un accouplement à valeurs dans $\mathbb{Z}/m\mathbb{Z}$
qui ne voit que ces morceaux. Pour une variété quelconque, il ramène la
question à la variété abélienne qui lui est attachée (son Albanese), par
une fonctorialité. Entre deux, une suite sur les surfaces : on découpe
une surface par un pinceau de courbes, et le nombre de classes de courbes
indépendantes sur la surface se lit sur le nombre de points rationnels
indépendants de la jacobienne de la courbe générique ; la question qui
clôt la suite — l'intersection sur la surface est-elle, au signe près, la
hauteur sur la jacobienne ? — est restée telle sur la page.

La dernière suite reprend tout depuis les fondations, à la manière de
Néron et Weil : pour un corps muni d'une famille de valeurs absolues, une
fonction « admissible » est une fonction de deux variables, un point et
une place, qui reste bornée là où les coordonnées le sont ; une
« quasi-fonction » est localement une fonction rationnelle corrigée par
une fonction admissible. Grothendieck en fait un faisceau, et la
théorie des hauteurs prend la forme d'une suite exacte de faisceaux :
toute classe de diviseurs porte une fonction « hauteur locale », unique à
une fonction bornée près. C'est le geste familier de ces années —
remplacer une collection d'estimations par une structure — appliqué à
l'arithmétique.

On a ici des notes de travail. Les questions sont écrites comme telles
(« À examiner », « à vérifier », « Cette condition est-elle
suffisante ?? »), les points admis sont dits admis, deux pages sont
barrées et une suite s'arrête au milieu d'une phrase. Le titre du dossier
est celui des archivistes ; la datation « [vers 1971] » est celle de
l'inventaire. Des pages dactylographiées d'autres textes y sont
mêlées ; elles ne sont pas lues ici.

Les noms sous lesquels chercher la suite : symbole local de Néron et
fonctions de Néron, modèle de Néron et groupe des composantes,
accouplement de Grothendieck sur les groupes de composantes, hauteur de
Néron–Tate ; formule de Shioda–Tate et réseaux de Mordell–Weil ; accouplement
de hauteur de Beilinson–Bloch ; fonctions de Weil, parties $M_K$-bornées et
« machine des hauteurs » de Weil.

\keywords{Néron local symbol, Néron model, component group, Grothendieck pairing, Néron-Tate height pairing, Picard scheme, Lefschetz pencil, Néron-Severi group, Shioda-Tate formula, Mordell-Weil lattice, K/k-trace, Tsen's theorem, Beilinson-Bloch height pairing, arithmetic intersection, Albanese variety, proper set of absolute values, MK-bounded set, Weil function, quasi-function, Weil height machine}
\end{resume}

\pagerange{1}{49}
\section*{Le fil du dossier, et les conventions}

\subsection*{Les stations}

Le dossier compte quarante-neuf feuillets, dont trente-quatre de sa main.
Ils se rangent en quatre suites, sans date, que rien ne relie sur la page
sinon le sujet\footnote{Le titre, « Hauteurs et symboles locaux », est
celui de l'inventaire, non le sien ; ses propres titres sont « Formules de
Néron » (feuillets 1 et 3) et « Quasi-fonctions de Néron-Weil » (feuillet
34). L'inventaire range le dossier dans le groupe « Descentes », daté
« [avant 1970] », alors que le dossier lui-même est daté « [vers 1971] » ;
on le constate sans le trancher.}.

\begin{itemize}
  \item[I.] \emph{Formules de Néron} (pages 1 à 7). Sur une variété
  abélienne $A$ définie sur le corps des fractions $K$ d'un anneau de
  valuation discrète, le symbole local de Néron est la somme d'un nombre
  d'intersection sur le modèle de Néron et d'un correctif $v(X,a)$ dont le
  dénominateur divise l'ordre $m$ du groupe des composantes ; d'où un
  accouplement $A^\vee(K) \times A(K) \to \mathbb{Z}/m\mathbb{Z}$. Puis
  (page 6) une construction du modèle de Néron d'une jacobienne comme
  adhérence dans le schéma de Picard.
  \item[II.] \emph{Une surface et un pinceau} (pages 9 à 14). Le rang du
  groupe de Néron–Severi d'une surface comparé au rang de Mordell–Weil de
  la jacobienne de la courbe générique d'un pinceau, et la question d'une
  formule de hauteur.
  \item[III.] \emph{Symboles locaux et nombres d'intersection} (pages 16 à
  31). Le programme : le symbole local de deux cycles algébriquement
  équivalents à zéro devrait être un nombre d'intersection sur un modèle
  régulier. La réduction, pour un diviseur et un $0$-cycle, au cas
  abélien de la suite I par l'application d'Albanese. La suite s'arrête
  sur le cas non lisse.
  \item[IV.] \emph{Quasi-fonctions de Néron–Weil} (pages 34 à 48). Les
  fondations : parties $M_K$-bornées, fonctions admissibles, faisceau des
  quasi-fonctions et ses suites exactes, sections globales.
\end{itemize}

Les suites I et III se répondent : la page 4 annonce en marge que les
formules « se généralise[nt] aussitôt aux variétés quelconques », et les
pages 25 à 31 sont cette généralisation, qui revient au cas abélien. La
suite II emprunte à la suite I la forme de Néron–Tate. La suite IV est
d'un autre registre — elle ne suppose aucun modèle, seulement des
valeurs absolues — et ne renvoie pas aux trois autres.

\subsection*{L'ordre des feuillets}

Dans la suite IV, l'ordre des archivistes n'est pas celui de la lecture :
la dernière phrase de la page 44 se poursuit en tête de la page 46, celle
de la page 46 en tête de la page 48, tandis que les pages 45 et 47, barrées
de longs traits, forment entre elles un fil distinct sur une autre
formulation de l'admissibilité. On lit ici $44 \to 46 \to 48$, puis
$45 \to 47$ à part ; rien sur la page ne dit lequel des deux fils a été
écrit le premier.

\subsection*{Conventions}

Les valuations sont écrites additivement : pour une valeur absolue
$|\ |_v$, $v(x) = -\log|x|_v$ à une constante positive près ; pour une
valuation discrète normalisée, $v(K^*) = \mathbb{Z}$. C'est la convention
des pages (feuillet 34 : « qu'on écrit sous forme logarithmique $v(x)$ »).

\begin{itemize}
  \item $\mathcal{A}$ désigne le modèle de Néron de $A$, $\mathcal{A}_0$
  sa fibre spéciale, $\mathcal{A}_0^0$ sa composante neutre,
  $\Phi_A = \mathcal{A}_0/\mathcal{A}_0^0$ le groupe des composantes et
  $m$ son ordre\footnote{Il écrit le modèle de Néron d'un $A$ à peine
  distinct de celui de la variété abélienne (page 3), puis d'un $A$ souligné
  (pages 23 et 31). La lettre $\Phi$ sert aussi, dans la suite IV, pour les
  fonctions admissibles ; le contexte les sépare, et l'on écrit $\Phi_A$
  avec indice pour le groupe des composantes.}.
  \item $A^\vee$ est la variété abélienne duale, qu'il note $A^*$.
  \item Les $0$-cycles sont notés $a, b$ (il les souligne). $D_a(A_K)$ et
  $D_\ell(A_K)$ sont les groupes de diviseurs algébriquement,
  linéairement équivalents à $0$ ; $Z_0^0(A_K)$ celui des $0$-cycles de
  degré $0$, et $Z_\ell(A_K)$ le noyau de l'application somme
  $Z_0^0(A_K) \to A(K)$.
  \item Le symbole local de la suite III est noté
  $\langle C_\eta, D_\eta\rangle_v$, comme aux pages 16 et 21 ; celui de la
  suite I $(X, a)_v$, comme à la page 4. Le second est le cas particulier
  du premier où la variété est abélienne et $C_\eta$ un $0$-cycle.
  \item Dans la suite IV, $\Phi_X$ est le faisceau des fonctions
  admissibles, $\mathcal{R}^*_X$ celui des fonctions rationnelles
  inversibles, $QF_X$ celui des quasi-fonctions, $\mathrm{Div}_X$ celui des
  diviseurs de Cartier, $U$ le groupe des $M_K$-unités de $K$\footnote{Il
  souligne les faisceaux ($\underline{\Phi}_X$, $\underline{QF}_X$,
  $\underline{R}^*_X$) ; les soulignements sont omis ici.}.
\end{itemize}

\subsection*{Ce que le dossier annonce sans l'établir}

Le corollaire de la page 5 est énoncé sans démonstration ; le « Lemme » de
la page 29 n'a ni énoncé lisible ni preuve ; l'indépendance du symbole
local vis-à-vis du modèle (page 16) est un programme, et les propriétés
de continuité et d'extension du corps de base (page 18) sont seulement
nommées. La formule de hauteur de la page 14 est posée comme question. Le
cas général b) de la page 31 s'arrête après deux lignes. Dans la suite IV,
le point de la page 41 est « admis », la condition de la page 45 reste une
question, et le cas de l'espace projectif (page 48), auquel tout est
ramené, n'est pas traité. Les notes de bas de page disent, cas par cas, ce
qui est vrai et ce que la littérature en sait.

\pagerange{1}{7}
\section*{I. Formules de Néron (pages 1 à 7)}

\pagerange{3}{5}
\subsection*{Le symbole de Néron sur une variété abélienne (pages 3 à 5)}

Soit $V$ un anneau de valuation discrète, de corps des fractions $K$, de
corps résiduel $k$, de valuation normalisée $v$\footnote{La page ajoute
« corps résiduel fini ? », avec un point d'interrogation. Avec $v$
normalisée par $v(K^*) = \mathbb{Z}$, rien de ce qui suit n'en dépend ; la
finitude du corps résiduel n'intervient que pour normaliser les valeurs
absolues dans le cas global de la page 5.}. Soit $A$ une variété abélienne
sur $K$, $\mathcal{A}$ son modèle de Néron sur $V$. On considère les couples
$(X, a)$ formés d'un diviseur $X \in D_a(A_K)$ et d'un $0$-cycle
$a \in Z_0^0(A_K)$, de supports disjoints\footnote{La page écrit
$Z_0^{*}(A)_K$ (l'exposant est douteux) sans préciser le degré. Que les
$0$-cycles soient de degré $0$ est notre lecture : elle est nécessaire pour
que l'accouplement de la page 4 se factorise par $A(K)$, via l'application
somme $Z_0^0(A_K) \to A(K)$, $\sum n_i (x_i) \mapsto \sum n_i x_i$. De même,
nous lisons $Z_\ell^{*}$ comme le noyau de cette application, qui est la
notion d'« équivalence linéaire » des $0$-cycles de degré $0$ sur une
variété abélienne.}.

Le feuillet 3 caractérise une fonction $v(X, a)$, à valeurs dans
$\mathbb{Q}$, définie sur ces couples, par quatre propriétés :

\begin{itemize}
  \item[(i)] elle est bilinéaire ;
  \item[(ii)] elle est invariante par les translations par les points de
  $A(K)$ ;
  \item[(iii)] si $X = \mathrm{div}(f)$, alors
  \[
    v(\mathrm{div}(f), a) = i\bigl(\mathrm{div}(\bar f)_{\mathrm{vert}},\ \bar a\bigr),
  \]
  où $\bar f$ est $f$ vue comme fonction rationnelle sur $\mathcal{A}$,
  $\mathrm{div}(\bar f)_{\mathrm{vert}}$ la partie verticale (portée par
  $\mathcal{A}_0$) de son diviseur, $\bar a$ l'adhérence de $a$ dans
  $\mathcal{A}$, et $i(\cdot,\cdot)$ la somme sur les points
  $x_0 \in \mathcal{A}_0$ des multiplicités d'intersection locales ;
  \item[(iv)] pour $X$ fixé, avec $0 \notin \mathrm{Supp}\,X$, la fonction
  $x \mapsto v(X, (x) - (0))$, définie sur les $x \in A(K)$ hors de
  $\mathrm{Supp}\,X$, est bornée.
\end{itemize}

Le modèle de Néron est lisse sur $V$, donc régulier : ses diviseurs de Weil
sont de Cartier, et les multiplicités d'intersection de (iii) ont un sens
dès que $\bar a$ ne rencontre pas la partie horizontale de
$\mathrm{div}(\bar f)$ hors de la fibre spéciale, ce qu'assure la
disjonction des supports\footnote{Pour un $0$-cycle dont les points ne
sont pas rationnels sur $K$, l'adhérence $\bar a$ est un $1$-cycle
quasi-fini sur $V$, non nécessairement fini puisque $\mathcal{A}$ n'est pas
propre ; les multiplicités locales aux points de $\mathcal{A}_0$ restent
définies, et en nombre fini. La page n'en dit rien.}.

Il en tire, sans démonstration écrite, les propriétés suivantes.

\begin{itemize}
  \item[(iv bis)] Pour $X$ et $a \in A(K)$ fixés, $v(X, (x) - (a))$ ne
  dépend que de la classe de $x$ modulo le sous-groupe $A^0(K)$ des points
  dont la réduction tombe dans $\mathcal{A}_0^0$. Comme $A(K)/A^0(K)$
  s'injecte dans le groupe fini $\Phi_A(k)$, elle ne prend qu'un nombre
  fini de valeurs\footnote{Tout point de $A(K)$ se prolonge en une section
  de $\mathcal{A}$ (propriété de Néron), d'où l'application de réduction
  $A(K) \to \mathcal{A}_0(k) \to \Phi_A(k)$, de noyau $A^0(K)$.}.
  \item[(v)] Si $X \in D_\ell(A_K)$ ou $a \in Z_\ell(A_K)$, alors
  $v(X, a) \in \mathbb{Z}$.
  \item[(vi)] En général, $m\, v(X, a) \in \mathbb{Z}$, où $m$ est l'ordre
  de $\Phi_A$\footnote{La page dit « $m =$ degré de $A_0/A_0^0$ », c'est-à-dire
  l'ordre du schéma en groupes fini étale $\Phi_A$, soit le cardinal de
  $\Phi_A(\bar k)$.}.
\end{itemize}

Une marge précise le cas extrême : si $m = 1$, c'est-à-dire si
$\mathcal{A}_0$ est connexe, alors $v$ est identiquement nulle\footnote{La
note marginale est en partie illisible ; ce qu'on en lit (« si
$X \sim_{\mathrm{alg}} 0$ … vraiment $v(X,a) = 0$ pour tous $X, a$ ») est
cohérent avec (iv bis) : si $\Phi_A = 0$, $v(X, (x)-(a))$ ne dépend pas de
$x$, donc est nulle en prenant $x = a$, et ces cycles engendrent
$Z_0^0(A_K)$.}.

\emph{Un accouplement à valeurs dans $\mathbb{Z}/m\mathbb{Z}$.} Par
(v) et (vi), $v$ induit, modulo $\mathbb{Z}$, un accouplement
\[
  D_a(A_K)/D_\ell(A_K) \times Z_0^0(A_K)/Z_\ell(A_K)
  \longrightarrow \tfrac{1}{m}\mathbb{Z}/\mathbb{Z} \simeq \mathbb{Z}/m\mathbb{Z},
\]
c'est-à-dire, puisque $D_a/D_\ell = \mathrm{Pic}^0(A_K) = A^\vee(K)$ et que
l'application somme identifie $Z_0^0/Z_\ell$ à $A(K)$,
\[
  A^\vee(K) \times A(K) \longrightarrow \mathbb{Z}/m\mathbb{Z}.
\]
Par (iv bis), il se factorise à droite par $A(K)/A^0(K) \subset \Phi_A(k)$.
La page le dit « remarquable » et note en marge, entourée : « À examiner :
cas d'une polarisation principale, jacobienne d'une variété
abélienne »\footnote{Lecture douteuse pour « polarisation principale »
et « jacobienne ». Un accouplement entre groupes de composantes
$\Phi_{A^\vee} \times \Phi_A \to \mathbb{Q}/\mathbb{Z}$ porte aujourd'hui son
nom : c'est l'accouplement de Grothendieck de SGA~7, exposé~IX, et son lien
avec le symbole de Néron a été étudié ensuite (Bosch et Lorenzini, entre
autres ; cité de mémoire). Ces feuillets ne citent pas SGA~7 et ne
définissent l'accouplement que par la fonction $v$ ; que les deux
coïncident, à un signe près, nous ne l'établissons pas ici.}.

\emph{La relation avec le symbole de Néron.} Soit $(X, a)_v$ le
symbole local de Néron, normalisé par $(\mathrm{div}\,f, a)_v = v(f(a))$,
où $f(a) = \prod f(x_i)^{n_i}$ pour $a = \sum n_i (x_i)$. La page encadre
la formule
\[
  \boxed{\ (X, a)_v = v(X, a) + i(\bar X, \bar a)\ }
  \qquad \text{(vii)}
\]
où $\bar X$ est l'adhérence de $X$ dans $\mathcal{A}$ et $i(\bar X, \bar a)$
la somme des multiplicités d'intersection aux points de la fibre spéciale.
Pour $X = \mathrm{div}(f)$ elle est compatible avec (iii) :
$\mathrm{div}(\bar f) = \bar X + \mathrm{div}(\bar f)_{\mathrm{vert}}$, et
$v(f(a)) = i(\mathrm{div}(\bar f), \bar a)$ est la somme des deux termes.
Comme $i(\bar X, \bar a)$ est entier, (vi) donne
$(X, a)_v \in \frac{1}{m}\mathbb{Z}$ ; en particulier le symbole de Néron
est \emph{rationnel}, et entier si $X \in D_\ell$ (on le savait :
c'est $v(f(a))$) ou si $a \in Z_\ell$. On obtient ainsi une autre
expression de l'accouplement de la page 4, en termes du seul symbole de
Néron. La marge ajoute : « se généralise aussitôt aux variétés
quelconques » ; c'est l'objet de la suite III.

\emph{Un corollaire de symétrie (page 5).} Pour un diviseur $X$ et un
$0$-cycle $a = \sum n_i(x_i)$, notons $X_a = \sum n_i\, \tau_{x_i}(X)$
la somme des translatés, $X^- = [-1]^* X$ et $a^- = [-1]_* a$\footnote{Ces
notations ne sont pas définies sur la page ; c'est notre lecture de
$X_{\underline a}$, $X^{-}$ et $\underline a^{-}$. Si $\deg a = 0$, le
théorème du carré fait de $X_a$ un diviseur algébriquement équivalent à
$0$, de classe $\varphi_X(\Sigma a)$ au signe près, où
$\varphi_X : A \to A^\vee$ est l'homomorphisme attaché à $X$.}. La page
énonce
\[
  v(X_a, b) = v(X^-_b, a) \qquad \bigl(= v(X_{b^-}, a^-)\bigr),
\]
sans démonstration. L'analogue global est vrai : l'homomorphisme
$\varphi_X$ est symétrique et ne change pas quand on remplace $X$ par
$X^-$, de sorte que la forme de Néron–Tate
$\langle \varphi_X(x), y\rangle$ est symétrique en $x, y$. L'énoncé local
n'est pas vérifié ici.

\emph{Le cas global (page 5).} Supposons donnée sur $K$ une famille de
valuations \emph{discrètes} $v$ satisfaisant la formule du produit
$\sum_v n_v\, v(x) = 0$ pour $x \in K^*$, avec des poids $n_v$ entiers
— c'est le cas d'un corps de fonctions d'une courbe, et ce n'est pas celui
d'un corps de nombres, qui a des places archimédiennes. La somme
$\sum_v n_v (X, a)_v$ ne dépend plus que des classes de $X$ modulo
$D_\ell$ et de $a$ modulo $Z_\ell$ — c'est la formule du produit pour
$X = \mathrm{div}\,f$ — et définit un accouplement
\[
  A^\vee(K) \times A(K) \longrightarrow \mathbb{R}
\]
\footnote{La page écrit $A^{*}(K) \times A^{*}(K)$, les deux astérisques
surchargés ; on attend $A^\vee(K) \times A(K)$, comme à la page 4.}, qui est
au signe près la forme de hauteur de Néron–Tate\footnote{La page dit
« la forme de Néron-Tate » sans signe. Avec des valuations additives, une
hauteur est grande là où les coordonnées ont des pôles, c'est-à-dire où
les valuations sont négatives : la somme des symboles locaux est
l'\emph{opposée} de la forme de Néron–Tate. C'est le signe qu'écrit la page
21 : $-\langle D_\eta, C_\eta\rangle = \sum_v \langle D_\eta,
C_\eta\rangle_v$.}. Il est à valeurs dans $\frac{1}{m}\mathbb{Z}$, où $m$
est le plus petit commun multiple des ordres $m_v$ des groupes de
composantes\footnote{La page dit « $m$ un entier $> 0$ convenable ». Cela
demande que les poids $n_v$ soient entiers et que $m_v = 1$ pour presque
tout $v$ (bonne réduction hors d'un ensemble fini), ce qui est le cas ; on
prend alors $m = \mathrm{ppcm}_v\, m_v$. Pour un corps de nombres,
l'énoncé est faux : les contributions archimédiennes sont des réels
quelconques.}.

\pagerange{6}{6}
\subsection*{Le modèle de Néron d'une jacobienne comme adhérence (page 6)}

Soit $f : X \to S$ un morphisme propre et plat, $X$ régulier, $S$ régulier
de dimension $1$, de point générique $\eta$, tel que
$\underline{\mathrm{Pic}}_{X/S}$ existe. Soit $\mathcal{A}$ l'adhérence
schématique de $A_\eta = \underline{\mathrm{Pic}}^0_{X_\eta/k(\eta)}$ dans
$\underline{\mathrm{Pic}}_{X/S}$. Supposons que $\mathcal{A}$ soit un
sous-schéma en groupes \emph{lisse} sur $S$ — ce qui a lieu dès que
$\underline{\mathrm{Pic}}_{X/S}$ est lisse sur $S$ le long de la section
unité, par exemple si les $\underline{\mathrm{Pic}}_{X_s/k(s)}$ sont
lisses — et \emph{séparé}, et que $X$ admette des sections étale-localement
sur $S$. Alors $\mathcal{A}$ est le modèle de Néron de $A_\eta$\footnote{Deux
hypothèses sont ajoutées. La séparation : quand la fibre spéciale est
réductible, $\underline{\mathrm{Pic}}_{X/S}$ n'est pas séparé, l'adhérence
de $A_\eta$ contient l'adhérence de la section unité, et un modèle de
Néron doit être séparé ; il faut alors quotienter par cette adhérence.
L'existence de sections étale-locales : sans elle, un point de
$\underline{\mathrm{Pic}}_{X_\eta}$ n'est pas toujours la classe d'un
faisceau inversible sur $X_\eta$ ; c'est l'avertissement de sa marge
(« Attention s'il n'y a pas de section locale (ét.) de $X$ sur $S$ »). La
page envisage $S$ régulier de dimension quelconque et ne restreint à la
dimension $1$ que l'exemple ; nous énonçons en dimension $1$. Sous ces
hypothèses, l'énoncé est aujourd'hui le théorème de Raynaud
(« Spécialisation du foncteur de Picard », 1970 ; Bosch–Lütkebohmert–Raynaud,
\emph{Néron Models}, chap.~9 ; cité de mémoire). La page ne cite personne.}.

La raison qu'il donne est la bonne : si $S' \to S$ est lisse,
$X' = X \times_S S'$ est lisse sur $X$ régulier, donc régulier ; sur un
schéma régulier tout diviseur de Weil est de Cartier, de sorte que tout
faisceau inversible sur la fibre générique $X'_{\eta'}$ se prolonge à $X'$
(on prend l'adhérence d'un diviseur). Tout point de $A_\eta$ à valeurs dans
la fibre générique de $S'$ se prolonge donc en un $S'$-point de
$\underline{\mathrm{Pic}}_{X/S}$, qui tombe dans $\mathcal{A}$ parce que
$S'$ est plat sur $S$ : c'est la propriété d'extension de Néron.

\pagerange{7}{7}
\subsection*{Une liste en marge (page 7)}

Le feuillet 7 est le sommaire d'une rédaction dactylographiée de Bourbaki
sur les systèmes de racines ; on n'en retient que cinq lignes au crayon
dans sa marge : « Matrices de Coxeter ; Matrices de Cartan ; Systèmes de
racines ; groupes engendrés par des réflexions ; diagrammes de
Dynkin »\footnote{La transcription note que rien n'assure que ces lignes
soient de sa main. Elles n'ont pas de rapport avec le reste du dossier.
Les feuillets 17 et 19 sont deux autres pages de la même rédaction, sans
rien de sa main.}.

\pagerange{9}{14}
\section*{II. Le groupe de Néron–Severi d'une surface et un pinceau (pages 9 à 14)}

\pagerange{9}{11}
\subsection*{Le pinceau et l'éclaté (pages 9 à 11)}

Soit $X$ une surface projective lisse et connexe sur un corps $k$
algébriquement clos\footnote{Une note entourée ajoute : « n'est pas une
restriction essentielle ! $\dim X \geqslant 2$ ». Tout ce qui suit est écrit
pour une surface, et la formule de la page 13 en dépend ; on s'y tient.},
plongée dans $\mathbb{P}^n$. Un pinceau de sections hyperplanes, défini par
un sous-espace linéaire $L \simeq \mathbb{P}^{n-2}$, coupe $X$
transversalement en un ensemble fini $Z$ ; soit $Y_0$ une section lisse du
pinceau, et $z = \mathrm{card}\, Z = (Y_0 \cdot Y_0)$. Soit
$\varphi : \tilde X \to X$ l'éclaté de $Z$ ; le pinceau devient un morphisme
$\tilde X \to \mathbb{P}^1$, de fibres les $\tilde Y_t \simeq Y_t$, de fibre
générique $Y_\eta$ sur $K = k(\eta) = k(t)$. Les courbes exceptionnelles
$\tilde z$ ($z \in Z$) sont des sections.

Notons $\mathrm{Pic}(X)'$ et $\mathrm{Pic}(\tilde X)'$ les sous-groupes des
classes dont la restriction à la fibre générique géométrique
$Y_{\bar\eta}$ est de degré $0$\footnote{La page écrit par endroits
$\mathrm{Pic}(X^*)'$, l'exposant douteux et absent de la définition, et
$\tilde X_{\bar\eta}$ pour $\tilde Y_{\bar\eta}$ ; on lit $\mathrm{Pic}(X)'$
et $Y_{\bar\eta}$.} ; pour $X$, ce sont les classes $\xi$ telles que
$\xi \cdot Y_0 = 0$. Il s'agit de déterminer les noyaux et les images des
restrictions
\[
  \mathrm{Pic}(X)' \hookrightarrow \mathrm{Pic}(\tilde X)'
  \longrightarrow \mathrm{Pic}^0(Y_{\bar\eta}) = J(Y_{\bar\eta}).
\]

\emph{Les noyaux.} La restriction
$\mathrm{Pic}(Y_\eta) \to \mathrm{Pic}(Y_{\bar\eta})$ est injective (la
courbe $Y_\eta$ a des points rationnels, les sections $\tilde z$). Le noyau
de $\mathrm{Pic}(\tilde X)' \to J(Y_{\bar\eta})$ est donc celui de
$\mathrm{Pic}(\tilde X) \to \mathrm{Pic}(Y_\eta)$ : les classes des diviseurs
\emph{verticaux}. Les fibres $\tilde X_y$, $y \in \mathbb{P}^1(k)$, sont
linéairement équivalentes entre elles. Si toutes les fibres sont
irréductibles, le noyau est donc engendré par la classe
\[
  \tilde\xi_0 = [\tilde Y_0] = \varphi^*(\xi_0) - \sum_{z \in Z} [\tilde z],
\]
où $\xi_0$ est la classe de $Y_0$ dans $X$\footnote{La page demande
d'abord les fibres « intègres », puis corrige page 13 : il suffit
qu'elles soient irréductibles, « on admet des sections multiples ». En fait,
un pinceau qui a des points base n'a pas de fibre multiple, puisque chaque
fibre rencontre une section $\tilde z$ avec multiplicité $1$ ;
« irréductible » et « intègre » reviennent donc ici au même.}. Comme
$\mathrm{Pic}(\tilde X) = \varphi^*\mathrm{Pic}(X) \oplus
\bigoplus_{z} \mathbb{Z}[\tilde z]$ et que $Z \neq \emptyset$, aucun multiple
non nul de $\tilde\xi_0$ ne provient de $X$ : la restriction
$\mathrm{Pic}(X)' \to \mathrm{Pic}(Y_{\bar\eta})$ est \emph{injective}.

\emph{La partie fixe.} L'image de
$\mathrm{Pic}^0(X) = \mathrm{Pic}^0(\tilde X)$ dans $J(Y_{\bar\eta})$ est
la « partie fixe » $J^{\mathrm{fixe}}$ de la jacobienne $J = J(Y_\eta)$ —
la $K/k$-trace, dans le vocabulaire d'aujourd'hui —, et
$J^{\mathrm{fixe}}(k) = J^{\mathrm{fixe}}(k(t))$ parce que $k(t)$ est une
extension transcendante pure de $k$ : une application rationnelle de
$\mathbb{P}^1$ dans une variété abélienne est constante\footnote{La page
écrit « extension pure de $k(\eta)$ », lecture douteuse ; le sens exige
« de $k$ ». Elle écrit aussi des \emph{isomorphismes}
$\mathrm{Pic}^0(X) \simeq \mathrm{Pic}^0(\tilde X) \simeq
J^{\mathrm{fixe}}(k)$. Le premier est vrai ; le second n'est assuré en
général qu'à isogénie près (le noyau de $\mathrm{Pic}^0(X) \to
\mathrm{Pic}^0(Y)$ est fini, nul par exemple en caractéristique $0$). Seuls
les rangs servent dans la suite, et l'énoncé à isogénie près suffit.}.
Posons $J' = J / J^{\mathrm{fixe}}$, la « partie mobile ». On obtient
\[
  NS(X)' \hookrightarrow NS(\tilde X)' \longrightarrow J'(K),
\]
où $NS(\cdot)'$ désigne l'image de $\mathrm{Pic}(\cdot)'$ dans le groupe de
Néron–Severi.

\pagerange{12}{13}
\subsection*{Le rang (pages 12 et 13)}

\emph{Le noyau.} Celui de $NS(\tilde X)' \to J'(K)$ est engendré par
l'image de $\tilde\xi_0$ : si $\alpha$ est dans le noyau, il provient
d'un $\alpha_0 \in \mathrm{Pic}(\tilde X)'$ dont l'image dans $J(K)$ est
fixe, donc provient de $\mathrm{Pic}^0(\tilde X)$ ; quitte à modifier
$\alpha_0$ par un élément de $\mathrm{Pic}^0(\tilde X)$, qui ne change pas
sa classe dans $NS$, son image est nulle et $\alpha_0$ est vertical. Il est
donc de rang $1$ (toujours sous l'hypothèse des fibres irréductibles).

\emph{L'image.} $NS(\tilde X)' \to J'(K)$ est surjectif modulo un
groupe fini. D'abord, $J(K) \to J'(K)$ l'est : son conoyau s'injecte dans
$H^1(K, J^{\mathrm{fixe}})$, groupe de torsion, et $J'(K)$ est de type fini
(théorème de Lang–Néron), d'où un conoyau fini. Ensuite
$\mathrm{Pic}(\tilde X)' \to J(K)$ est surjectif : un point de $J(K)$ est
la classe d'un diviseur de degré $0$ sur $Y_\eta$ — le groupe de Brauer de
$k(t)$ est nul, par le théorème de Tsen — et un diviseur sur $Y_\eta$ se
prolonge à $\tilde X$ par adhérence\footnote{La première de ces deux
étapes est écrite sur la page, sans le théorème de Lang–Néron, qui
n'est pas cité ; nous l'ajoutons, car « surjectif modulo un groupe de
torsion » ne suffirait pas pour une égalité de rangs si $J'(K)$
n'était pas de type fini. Le point rationnel des sections $\tilde z$
dispense d'ailleurs de Tsen.}.

\emph{La formule.} Comme $NS(\tilde X) \to \mathbb{Z}$,
$\xi \mapsto \xi \cdot F$ ($F$ la classe d'une fibre) est surjective (une
section $\tilde z$ est de degré $1$ sur les fibres), on a
$\rho(\tilde X) = 1 + \mathrm{rg}\, NS(\tilde X)' = 2 + \mathrm{rg}\, J'(K)$,
et $\rho(\tilde X) = \rho(X) + z$. La page encadre :
\[
  \boxed{\ \rho(\tilde X) - 2 = \rho(X) + z - 2 = \rho(J')\ }
\]
où $\rho$ est le nombre de Picard, $\rho(J')$ le rang de $J'(K)$, et $z$ le
nombre de points base du pinceau\footnote{La page définit $z$ comme le
« nb des sections singulières du pinceau », lecture marquée douteuse. Avec
ce sens la formule serait fausse : c'est l'éclatement des $z$ points base
qui ajoute $z$ au nombre de Picard ; $z = \mathrm{card}\,Z = (Y_0 \cdot
Y_0)$, le degré de $X$. Les deux « 2 » sont repassés sur la page.},
valable quand toutes les fibres sont irréductibles. C'est, pour ce
pinceau, la forme que prend la formule dite aujourd'hui de Shioda–Tate,
qui dans le cas général ajoute au second membre la contribution
$\sum_t (r_t - 1)$ des fibres réductibles à $r_t$
composantes\footnote{Shioda, « On the Mordell–Weil lattices », 1990, et
Tate pour les surfaces elliptiques ; cité de mémoire. Les feuillets ne
citent pas ces noms.}.

\pagerange{14}{14}
\subsection*{Une formule de hauteur posée en question (page 14)}

La jacobienne $J'$ hérite de $J$ une polarisation : la duale de $J'$
s'envoie dans la duale de $J$, qui est $J$ munie de sa polarisation
principale, et l'on restreint\footnote{Cette explication est biffée sur la
page, sans être remplacée ; on la garde parce qu'elle est juste.}. Par la
théorie de Néron–Tate, cette polarisation définit sur $J'(K)$ une forme
quadratique $Q$, définie positive sur $J'(K) \otimes \mathbb{R}$ ($J'$ est
de $K/k$-trace nulle). D'autre part, la forme d'intersection est définie
\emph{négative} sur $NS(X)' \otimes \mathbb{R}$, orthogonal de la classe
ample $\xi_0$ : c'est le théorème de l'indice de Hodge. Soit
$i : NS(X)' \to J'(K)$ l'application précédente. La page propose, avec un
point d'interrogation :
\[
  \boxed{\ Q(i(\xi), i(\eta)) = -\,\xi \cdot \eta\ } \qquad
  (\xi, \eta \in NS(X)'),
\]
« valable si les $Y_t$ sont tous intègres ? ».

La réponse est oui, au choix près de la normalisation de $Q$ : c'est la
formule de hauteur de Shioda\footnote{Réponse nôtre. Dans la théorie des
réseaux de Mordell–Weil, la hauteur d'un point $P$ est $-(\psi(P) \cdot
\psi(P))$, où $\psi(P)$ est la projection de la classe de $\bar P - \bar O$
orthogonalement au réseau engendré par la section nulle, la fibre et les
composantes des fibres ne rencontrant pas la section nulle. Ici, toutes
les fibres étant irréductibles, ce réseau est engendré par une section
$\tilde z_0$ et par $F$ ; or $\varphi^*\xi$ est orthogonal à $F$ (car
$\xi \cdot Y_0 = 0$) et aux $\tilde z$, et ne diffère de la classe
$\bar P - \bar O$ du point $P = i(\xi)$ que par un multiple de $F$ : on a
donc $\psi(i(\xi)) = \varphi^*\xi$ et
$\langle i\xi, i\eta\rangle = -\varphi^*\xi \cdot \varphi^*\eta = -\xi\cdot\eta$.
La hauteur de Shioda est la forme de Néron–Tate de $J$ pour sa
polarisation principale, à un facteur de normalisation près, et elle est
nulle sur la partie fixe, donc se lit sur $J'(K)$ ; que la polarisation
« héritée » de $J'$ donne exactement ce facteur, nous ne le vérifions pas.
Sans l'hypothèse des fibres irréductibles, $\varphi^*\xi$ n'est plus
orthogonal aux composantes des fibres et la formule demande une
correction : la condition de la page est la bonne.}.

\pagerange{16}{31}
\section*{III. Symboles locaux et nombres d'intersection (pages 16 à 31)}

\pagerange{16}{18}
\subsection*{Le programme (pages 16 et 18)}

Soit $f : X \to S$ un morphisme projectif, $X$ régulier, $S$ le spectre
d'un anneau de valuation discrète de point générique $\eta$ et de point
spécial $s$ ; supposons $X_\eta$ lisse de dimension $n$ et $X_s$
géométriquement intègre. Soient $C$, $D$ deux cycles sur $X$, de
codimensions $p$ et $q$ avec $p + q = n + 1$\footnote{La page écrit d'abord
« de dimensions $c$ et $d$ avec $c + d = n + 1$ », puis « de codimensions
$c, d$ tels que $c + d = n + 1$ » pour les fibres génériques. Les deux
sommes valent $n + 1$ — pour des cycles horizontaux, codimension dans $X$
et codimension de la fibre générique dans $X_\eta$ coïncident —, mais les
indices ne désignent pas la même chose ; on indexe par la codimension.}.
On suppose :

\begin{itemize}
  \item[a)] $C_\eta$ et $D_\eta$ sont algébriquement équivalents à $0$ sur
  $X_\eta$ ;
  \item[b)] $C$ et $D$ n'ont que des composantes horizontales :
  $C = \overline{C_\eta}$, $D = \overline{D_\eta}$ ;
  \item[c)] l'intersection $C \cdot D$ est propre, ce qui force
  $\mathrm{Supp}\,C_\eta \cap \mathrm{Supp}\,D_\eta = \emptyset$.
\end{itemize}

Alors $C \cdot D$ est un $0$-cycle porté par $X_s$, de degré $i(C \cdot D)$,
et ce degré ne dépend que de $C_\eta$ et $D_\eta$. Le programme est de le
décrire directement en termes de $C_\eta$ et $D_\eta$, comme un symbole
\[
  \langle C_\eta, D_\eta\rangle_v \in \mathbb{Q}
\]
\emph{indépendant du modèle $X$}, et même de son existence\footnote{« à
priori », dans une bulle, se rattache à « à valeurs dans $\mathbb{Q}$ ».
L'indépendance vis-à-vis du modèle est ce que le programme cherche, non
un acquis. C'est aujourd'hui l'accouplement de hauteur locale de
Beilinson et Bloch (années 1980 ; cité de mémoire) ; au-delà des diviseurs
il demande en général de corriger $\overline{C_\eta}$ par un cycle vertical
qui le rende orthogonal aux cycles portés par la fibre spéciale. Quand
$X$ est lisse sur $S$, cette correction est inutile, la spécialisation d'un
cycle homologiquement trivial l'étant encore.}.

Plus généralement (page 18), pour des cycles $Z_1, \ldots, Z_k$ de
codimensions $c_1 + \cdots + c_k = n + 1$, dont deux ont une fibre générique
algébriquement équivalente à $0$, et dont l'intersection est propre, on
aurait $i(Z_1 \cdots Z_k) = \langle Z_{1\eta}, \ldots, Z_{k\eta}\rangle_v$.
Les propriétés attendues du symbole local sont :

\begin{itemize}
  \item[(i)] multilinéarité et symétrie ;
  \item[(ii)] invariance par les isomorphismes de $X_\eta$ ;
  \item[(ii bis)] pour $k = 2$, une formule de transposition pour un morphisme
  $g : X_\eta \to Y_\eta$, soit
  $\langle C, g^*D\rangle_v = \langle g_*C, D\rangle_v$ (formule de
  projection) ;
  \item[(iii)] normalisation : si $Z_{k\eta} = Z'_{k\eta} \cdot
  \mathrm{div}(f_\eta)$ et que $Z_{1\eta} \cdots Z_{k-1,\eta} \cdot
  Z'_{k\eta}$ est un $0$-cycle bien défini, alors
  \[
    \langle Z_{1\eta}, \ldots, Z_{k\eta}\rangle_v =
    v\bigl(f_\eta(Z_{1\eta} \cdots Z_{k-1,\eta} \cdot Z'_{k\eta})\bigr) ;
  \]
  \item[(iv)] un principe de continuité ;
  \item[(v)] la compatibilité à l'extension du corps de base.
\end{itemize}

Les deux dernières ne sont que nommées\footnote{La page commence par
« Le cas gl. », où « gl. » peut se lire « global » ou « général » ; le
contexte (une seule valuation) fait lire « général ». La forme explicite
de (ii bis) est nôtre : la page n'écrit que le nom. « (v) Extension du corps
de base » est écrit en marge.}.

\pagerange{21}{21}
\subsection*{L'énoncé (page 21)}

Dans la situation globale — $f : X \to S$ projectif, à fibres
géométriquement irréductibles, de fibre générique lisse, $X$ régulier, $S$
muni d'un ensemble $M$ de valuations discrètes satisfaisant la formule du
produit —, la page écrit, pour $C_\eta$ et $D_\eta$ algébriquement
équivalents à $0$,
\[
  -\langle D_\eta, C_\eta\rangle = \sum_{v \in M} \langle D_\eta,
  C_\eta\rangle_v = \sum_{v \in M} (D \cdot C)_v,
\]
où $\langle\ ,\ \rangle$ sans indice est l'accouplement de hauteur global.
Le signe est celui qu'on attend avec des valuations additives (voir la
page 5).

La situation locale est celle-ci : $S$ local complet, de corps résiduel
algébriquement clos, $f$ projectif, $X_\eta$ lisse, $X_0$ lisse sauf en un
nombre fini de points où $X$ est régulier ; $C$, $D$ sans composante
verticale, se coupant proprement (leurs supports se rencontrent dans
$X_0$) ; $C_\eta$, $D_\eta$ algébriquement équivalents à $0$. L'énoncé
visé est
\[
  \boxed{\ \langle C_\eta, D_\eta\rangle_v = i(C \cdot D)\ }
\]
\footnote{Plusieurs mots sont illisibles dans les hypothèses (la nature des
points singuliers de $X_0$, trois mots soulignés au bas du feuillet), et
« régulier en ces points » est une lecture douteuse. Le reste de la suite
ne traite que le cas où $D$ est un diviseur et $C$ un $0$-cycle sur la
fibre générique, c'est-à-dire le symbole de Néron de la suite I étendu à
une variété quelconque.}

\pagerange{23}{23}
\subsection*{Un brouillon sur les modèles de Néron (page 23)}

Le feuillet 23 n'a pas de phrase suivie. En tête, encadré : un anneau local
noethérien réduit $A$ de dimension $2$, de profondeur $2$, et
$f \in \mathfrak{m}_A$ tel que $A/fA$ soit régulier. C'est le critère
usuel : si $f$ n'est pas diviseur de zéro, $A/fA$ régulier entraîne $A$
régulier ; appliqué à une uniformisante, il dit que $X$ est régulier là où
la fibre spéciale l'est\footnote{L'usage de l'encadré n'est pas écrit ; le
rapprochement avec les hypothèses de la page 21 est nôtre.}. Le reste
consiste en éléments épars : une variété abélienne produit
$A_\eta = A_{1\eta} \times A_{2\eta}$ et les adhérences de ses facteurs
dans le modèle de Néron $\mathcal{A}$ ; un tore $T$ et une variété
abélienne $B$ ; des isogénies $B_\eta \xrightarrow{\alpha} A_\eta
\xrightarrow{\beta} B_\eta$ de composé $N$, et leurs prolongements aux
modèles de Néron ; une suite $0 \to A \to B \to C$. Aucun énoncé n'en
est tiré, et l'on n'en reconstitue pas.

\pagerange{25}{29}
\subsection*{La réduction au cas abélien (pages 25 à 29)}

\emph{La normalisation.} Pour $D = \mathrm{div}(f)$ sans composante
verticale, la propriété (iii) de la page 18 et le calcul direct donnent
\[
  \langle C_\eta, \mathrm{div}\,f_\eta\rangle_v = i(C \cdot \mathrm{div}\,f)
  = v\bigl(f_\eta(C_\eta)\bigr) ;
\]
la seconde égalité vaut pour tout $1$-cycle horizontal $C$, la longueur de
$\mathcal{O}_C/(f)$ en la fibre spéciale étant la valuation de la norme de
$f$ restreinte à $C$. Le feuillet 25 esquisse, pour une courbe plongée
dans sa jacobienne $J$, comment représenter un point $x \in X_\eta(K)$
comme diviseur d'une section rationnelle d'un fibré en droites $L$ sur $X$,
définie partout sauf au point critique $x_0$ et dont le diviseur n'a pas de
composante verticale\footnote{Le feuillet est très lacunaire (plusieurs
mots illisibles, deux croquis de fibres) ; on n'en retient que ce qui se
lit. Les flèches $X' \hookrightarrow \widetilde{X''} \to X'' \hookrightarrow
J$, avec $X''_0 = X_0$ et une flèche marquée « birat. », suggèrent une
modification birationnelle du modèle ; la page ne dit pas laquelle.}.

\emph{a) Le cas lisse (pages 27 et 29).} Supposons $f$ lisse. Alors
$X_\eta$ a une variété d'Albanese $A_\eta$ et une variété de Picard
$A^\vee_\eta$ ; soit $\mathcal{A}$ le modèle de Néron de $A_\eta$.
L'application d'Albanese $X_\eta \to A_\eta$ (définie par un point base) se
prolonge en $\alpha : X \to \mathcal{A}$, par la propriété de Néron puisque
$X$ est lisse sur $S$\footnote{La page renvoie au « th. de Weil » : une
application rationnelle d'un schéma lisse dans un schéma en groupes, définie
en codimension $1$, l'est partout. C'est ce théorème qui fonde la propriété
d'extension de Néron.}. Un diviseur $D_\eta$ algébriquement équivalent à $0$
est linéairement équivalent à $\alpha^*(D'_\eta)$ pour un diviseur $D'_\eta$
algébriquement équivalent à $0$ sur $A_\eta$.

Supposons d'abord $D = \alpha^*(D')$, où $D'$ est un diviseur sur
$\mathcal{A}$ tel que $\alpha^*(D')$ soit défini. Posons
$C' = \alpha_*(C)$, défini parce que $C$ est propre sur $S$. La formule de
projection donne
\[
  C' \cdot D' = \alpha_*(C) \cdot D' = \alpha_*\bigl(C \cdot \alpha^*D'\bigr)
  = \alpha_*(C \cdot D),
  \qquad\text{d'où}\qquad i(C \cdot D) = i(C' \cdot D'),
\]
et, par (ii bis), $\langle C_\eta, D_\eta\rangle_v = \langle C'_\eta,
D'_\eta\rangle_v$. La formule se réduit alors à
\[
  \boxed{\ i\bigl(\overline{C'_\eta} \cdot \overline{D'_\eta}\bigr) =
  \langle C'_\eta, D'_\eta\rangle_v\ }
\]
sur la variété abélienne, qui est la formule (vii) de la page 4 sans terme
correctif\footnote{Ce rapprochement est nôtre. $X$ étant lisse et propre
sur $S$, à fibres géométriquement intègres, l'Albanese a bonne réduction
— nous l'admettons, la page ne le discute pas —, de sorte que
$\mathcal{A}_0$ est connexe, $m = 1$, et $v(X, a) = 0$ par la marge de la
page 3.}.

Dans le cas général, on choisit (« Lemme », illisible\footnote{Le crochet
« [Lemme : …] » de la page 29 n'a ni énoncé lisible ni démonstration. Que
l'on puisse choisir $D'_\eta$ dans sa classe de façon que $D'_\eta$ ne
contienne pas $\alpha_\eta(X_\eta)$ et que son adhérence ne contienne pas
$\alpha_0(X_0)$ est un argument de position générale par translation ;
nous ne le détaillons pas.}) $D'_\eta$ tel que $D_\eta \sim
\alpha^*(D'_\eta)$, que $D'_\eta \not\supset \alpha_\eta(X_\eta)$ et que
$\overline{D'_\eta} \not\supset \alpha_0(X_0)$. Alors
\[
  D = \alpha^*(D') + \mathrm{div}(f),
\]
où $\alpha^*(D')$ est défini et sans composante verticale et $f$ est une
fonction rationnelle sur $X$. Comme $D$ et $\alpha^*(D')$ sont horizontaux,
$\mathrm{div}(f)$ est sans composante verticale, et la normalisation
conclut :
\[
  \boxed{\ \langle C_\eta, \mathrm{div}\,f_\eta\rangle_v = i(C \cdot
  \mathrm{div}\,f)\ } = v\bigl(f_\eta(C_\eta)\bigr)
  \qquad \text{si $\mathrm{div}(f)$ est sans composante verticale}
\]
\footnote{La page écrit « valable si $\mathrm{div}(f)$ est sans composante
\emph{horizontale} ». C'est un lapsus : $\mathrm{div}(f_\eta) = D_\eta -
\alpha^*(D'_\eta)$ est non nul en général, donc a des composantes
horizontales, et la condition dont l'argument a besoin — celle que la page
vient d'établir pour $\alpha^*(D')$ — est l'absence de composante
verticale ; c'est elle qui fait de $\mathrm{div}(f)$ l'adhérence de
$\mathrm{div}(f_\eta)$.}.

\pagerange{31}{31}
\subsection*{b) Le cas général, interrompu (page 31)}

Sans hypothèse de lissité, soit $\mathcal{A}$ le modèle de Néron de
$\mathrm{Alb}(X_\eta)$. L'ouvert $X' = X - \{x_0\}$, complémentaire du
point singulier de la fibre spéciale, est lisse sur $S$, et la propriété de
Néron fournit $\alpha : X' \to \mathcal{A}$\footnote{La page écrit « le pt
singulier de $X'$ », lecture douteuse ; le sens, et les hypothèses de la
page 21, exigent « de $X_0$ ». S'il y en a plusieurs, on les retire tous.}.
Le feuillet s'arrête là. La difficulté que la suite aurait eu à traiter
est visible : $\alpha$ n'est pas définie en $x_0$, et le modèle de Néron
n'a plus de raison d'avoir une fibre spéciale connexe, de sorte que le
terme correctif $v(X, a)$ de la page 4 doit reparaître.

\pagerange{34}{48}
\section*{IV. Quasi-fonctions de Néron–Weil (pages 34 à 48)}

Le titre est le sien\footnote{Le mot « quasi-fonction » est celui de
Néron (« Quasi-fonctions et hauteurs sur les variétés abéliennes »,
\emph{Ann. of Math.}, 1965 ; cité de mémoire). Le cadre des pages qui
suivent — ensemble propre de valeurs absolues, $M_K$-diviseurs, parties
$M_K$-bornées — est celui que Lang a exposé (\emph{Diophantine Geometry},
1962 ; \emph{Fundamentals of Diophantine Geometry}, 1983, chap.~10), où les
quasi-fonctions s'appellent « fonctions de Weil ». Les pages ne citent que
les deux noms du titre.}.

\pagerange{34}{36}
\subsection*{Fonctions sur $X(\bar K) \times M_{\bar K}$ (pages 34 à 36)}

Soit $K$ un corps muni d'un ensemble « propre » $M_K$ de valeurs absolues
(tout $x \in K^*$ est de valeur absolue $1$ pour presque toute
$v \in M_K$), écrites additivement. Soit $\bar K$ une clôture algébrique,
$M_{\bar K}$ l'ensemble des valeurs absolues de $\bar K$ prolongeant celles
de $M_K$, et $\Gamma = \mathrm{Aut}(\bar K/K)$. Les prolongements d'une
valeur absolue étant conjugués,
\[
  M_{\bar K}/\Gamma \xrightarrow{\ \sim\ } M_K. \qquad (*)
\]

Soit $X$ un schéma localement de type fini sur $K$, et
$\varphi : X(\bar K) \times M_{\bar K} \to \mathbb{R}$, $(P, v) \mapsto
\varphi(P, v)$, une fonction invariante par $\Gamma$. Si $P$ est fixé par
$\Gamma$, c'est-à-dire rationnel sur la clôture parfaite $K^{p^{-\infty}}$,
$\varphi(P, v)$ ne dépend que de $v|K$, par $(*)$ ; $\varphi$ induit donc
une fonction sur $X(K) \times M_K$. De même, pour toute extension finie $L$
de $K$ dans $\bar K$, elle induit $\varphi_L : X(L) \times M_L \to
\mathbb{R}$, avec la transitivité $\varphi_L(P, v|L) = \varphi_{L'}(P, v)$
pour $L \subset L'$, $P \in X(L)$, $v \in M_{L'}$ ; et la donnée de
$\varphi$ équivaut à celle d'un tel système compatible. On note
$\widehat\Phi(X) = \widehat\Phi(X/K, M_K)$ l'ensemble de ces $\varphi$ ; il
ne dépend pas, à isomorphisme unique près, du choix de $\bar K$. Pour $K'/K$
finie, on a $\widehat\Phi(X/K) \to \widehat\Phi(X_{K'}/K')$, bijective si
$K'/K$ est radicielle, et identifiant $\widehat\Phi(X/K)$ aux invariants
$\widehat\Phi(X_{K'}/K')^{\pi}$ si $K'/K$ est quasi-galoisienne de groupe
$\pi$\footnote{Les deux énoncés tiennent à ce que $\widehat\Phi(X_{K'}/K')$
est l'ensemble des fonctions invariantes par $\mathrm{Aut}(\bar K/K')$, qui
est égal à $\Gamma$ si $K'/K$ est radicielle et distingué de quotient $\pi$
si $K'/K$ est quasi-galoisienne.}.

\pagerange{36}{38}
\subsection*{Parties $M_K$-bornées et fonctions admissibles (pages 36 à 38)}

Un \emph{$M_K$-diviseur} est une fonction $c : M_K \to \mathbb{R}$ nulle
hors d'un ensemble fini. Soit $E$ un ensemble muni de $\pi : E \to
M_{\bar K}$ ; une fonction $\varphi$ sur $E$ est \emph{$M_K$-bornée} s'il
existe un $M_K$-diviseur $c$ avec $|\varphi(\xi)| \leqslant c(\pi(\xi)|K)$,
c'est-à-dire si $\varphi$ est bornée et nulle hors de $\pi^{-1}$ d'un
ensemble fini $S \subset M_K$.

Une partie $E$ de $X(\bar K) \times M_{\bar K}$ est \emph{$M_K$-bornée} si
elle est contenue dans une réunion finie d'ensembles
\[
  E(U; x_1, \ldots, x_n; c) = \Bigl\{(P, v) :\ P \in U(\bar K),\ \
  \min_{1 \leqslant i \leqslant n} v(x_i(P)) \geqslant c(v|K)\Bigr\},
\]
où $U$ est un ouvert affine de $X$, $x_1, \ldots, x_n$ des générateurs de
sa $K$-algèbre et $c$ un $M_K$-diviseur. Concrètement : il existe $S \subset
M_K$ fini et $M > 0$ tels que $P$ soit entier pour les $x_i$ en $v$ si
$v|K \notin S$, et que $v(x_i(P)) \geqslant -M$ si $v|K \in S$. Une marge
note que la notion est invariante par extension finie du corps de base,
que sa fonctorialité en $X$ « n'est pas claire », et pose une
question\footnote{« Dans le cas “global”, y a-t-il une notion “$M_K$-bornée”
de $\mathbf{P}^r_K$, celles des hauteurs bornées ? » ; quelques mots
illisibles. Pour un morphisme d'affines, l'image d'une partie $M_K$-bornée
l'est encore — les $x_i$ du but sont des polynômes en ceux de la source —,
ce qui donne la fonctorialité utilisée aux pages 39 et 48 ; remarque
nôtre. Sur la question, notons seulement que $\mathbf{P}^r(\bar K) \times
M_{\bar K}$ tout entier est $M_K$-borné (en chaque $(P, v)$, une
coordonnée de valuation minimale fournit une carte où toutes les autres
sont entières) : les parties de hauteur bornée relèvent d'une notion
globale que celle-ci ne voit pas. La page laisse la question ouverte.}.

Une fonction $\varphi$ est \emph{admissible} si elle est $M_K$-bornée sur
toute partie $M_K$-bornée. Les ensembles $\{P\} \times M_{\bar K}$ étant
$M_K$-bornés, pour $P$ fixé, $\tilde\varphi(P) = (v \mapsto \varphi(P,v))$
est majoré en module par un $M_K$-diviseur positif ; si $P$ est rationnel
sur une extension finie $K'$ et $\varphi$ invariante par $\Gamma$,
$\tilde\varphi(P)$ se factorise par $M_{K'}$ et est un $M_{K'}$-diviseur.
Avec $\mathrm{Div}(M_{\bar K}) = \varinjlim_{K'} \mathrm{Div}(M_{K'})$, on
obtient
\[
  \tilde\varphi : X(\bar K) \longrightarrow \mathrm{Div}(M_{\bar K}),
\]
compatible à l'action de $\Gamma$, et qui envoie toute partie $F$ de
$X(\bar K)$ telle que $F \times M_{\bar K}$ soit $M_K$-bornée sur une
partie bornée (pour l'ordre) de $\mathrm{Div}(M_{\bar K})$.

\pagerange{39}{41}
\subsection*{L'exemple fondamental, et le faisceau des fonctions admissibles (pages 39 à 41)}

Pour $f \in \Gamma(X, \mathcal{O}_X^*)$, la fonction
\[
  \varphi_f(P, v) = v\bigl(f(P)\bigr)
\]
est invariante par $\Gamma$ et admissible. Par fonctorialité ($f$ est un
morphisme $X \to \mathbf{G}_{m,K}$), on se ramène à $X = \mathbf{G}_{m,K}$
et $f = x$ la coordonnée\footnote{Une première rédaction, directe sur un
affine, est encadrée et barrée sur la page ; la fonctorialité invoquée est
celle de la note précédente.}. Une partie $M_K$-bornée de $\bar K^* \times
M_{\bar K}$ est contenue dans un ensemble $\{(x, v) : v(x) = 0$ si $v|K
\notin S$, $|v(x)| \leqslant M$ si $v|K \in S\}$, sur lequel $\varphi_x(x, v)
= v(x)$ est évidemment $M_K$-bornée.

Pour $U$ ouvert de $X$, soit $\Phi_X(U)$ l'ensemble des fonctions admissibles
$\Gamma$-invariantes sur $U(\bar K) \times M_{\bar K}$. C'est un préfaisceau
séparé ; pour qu'il soit un faisceau, il faut voir qu'une $\varphi$ dont les
restrictions aux $U_i$ d'un recouvrement ouvert sont admissibles l'est. On
se ramène à $X$ affine, et il suffit que toute partie $M_K$-bornée de
$X(\bar K) \times M_{\bar K}$ soit contenue dans une réunion finie d'images
de parties $M_K$-bornées des $U_i(\bar K) \times M_{\bar K}$. La page écrit
« à vérifier » et « Admettons ce point ». Il est vrai\footnote{Argument
nôtre. On peut supposer les $U_i = D(g_i)$ principaux, en nombre fini, et
écrire $1 = \sum h_i g_i$. Sur une partie bornée, $v(h_i(P)) \geqslant c(v)$
pour un $M_K$-diviseur $c$ ; comme $0 = v(1) \geqslant \min_i v(h_i g_i(P))
- c_0(v)$, où $c_0$ tient compte des places archimédiennes, il existe pour
chaque $(P, v)$ un $i$ tel que $v(g_i(P)) \leqslant c_1(v)$, c'est-à-dire
$v(1/g_i(P)) \geqslant -c_1(v)$ : $(P, v)$ est dans une partie bornée de
$D(g_i)$ pour les générateurs $x_1, \ldots, x_n, 1/g_i$. C'est, sous une
forme ou une autre, le lemme qui rend chez Lang la notion de partie bornée
locale sur $X$.}. On l'admet donc avec lui.

\pagerange{41}{44}
\subsection*{Le faisceau des quasi-fonctions (pages 41 à 44)}

On a deux homomorphismes de faisceaux, $\mathcal{O}_X^* \to \Phi_X$,
$f \mapsto \varphi_f$, et $\mathcal{O}_X^* \hookrightarrow \mathcal{R}^*_X$
(fonctions rationnelles inversibles ; $X$ réduit, voire séparable sur $K$,
« pour fixer les idées »). Pour obtenir « une généralisation commune des
fonctions admissibles et des fonctions rationnelles », on pose
\[
  QF_X = \Phi_X \amalg_{\mathcal{O}_X^*} \mathcal{R}^*_X
  = \bigl(\Phi_X \times \mathcal{R}^*_X\bigr)/\theta(\mathcal{O}_X^*),
  \qquad \theta(f) = (-\varphi_f,\ f),
\]
la somme amalgamée\footnote{L'exposant de $f$ dans $\theta(f)$ est un pâté
d'encre. Avec $\theta(f) = (-\varphi_f, f)$, la classe de $(\varphi, g)$ est
la fonction $(P, v) \mapsto \varphi(P, v) + v(g(P))$, bien définie là où $g$
l'est : $(\varphi - \varphi_u, gu)$ donne la même fonction. Le choix
$f^{-1}$ donnerait la convention de signe opposée ; c'est une convention,
et nous prenons celle-ci.}. Une quasi-fonction est donc, localement, une
fonction admissible plus la valuation d'une fonction rationnelle : une
fonction de Weil, dans le langage de Lang. On a deux suites exactes
\[
  \Phi_X \longrightarrow QF_X \xrightarrow{\ \mathrm{div}\ }
  \mathrm{Div}_X \longrightarrow 0, \qquad
  \mathcal{R}^*_X \longrightarrow QF_X \longrightarrow
  \Phi_X/\mathrm{Im}\,\mathcal{O}_X^* \longrightarrow 0,
\]
où $\mathrm{Div}_X = \mathcal{R}^*_X/\mathcal{O}_X^*$ est le faisceau des
diviseurs de Cartier.

\emph{$\Phi_X \to QF_X$ est injectif.} Si $(\varphi, 1) = (-\varphi_f, f)$,
alors $f = 1$, donc $\varphi = 0$\footnote{La page écrit l'égalité avec les
composantes dans l'ordre inverse, $(1, \varphi) = (f, -\varphi_f)$ ; c'est
sans conséquence.}. Autrement dit : \emph{les quasi-fonctions de diviseur
nul sont les fonctions admissibles.}

\emph{Le noyau de $\mathcal{R}^*_X \to QF_X$.} Il est formé des $f$
pour lesquels $(0, f) = (-\varphi_g, g)$ avec $g$ inversible, soit
$f = g \in \Gamma(X, \mathcal{O}_X^*)$ et $v(g(P)) = 0$ pour tous $P$ et $v$.
Si $g$ n'était pas constante sur une composante irréductible géométrique,
$g(X(\bar K))$ contiendrait le complémentaire d'une partie finie de
$\bar K^*$, et l'on trouverait des valeurs de valuation non nulle ; $g$ est
donc localement constante sur les composantes géométriques, algébrique sur
$K$, et partout une unité\footnote{La prose de la page 43 est d'une écriture
très rapide et en partie illisible ; seules les formules et la conclusion
se lisent sûrement. L'argument suppose qu'au moins une $v \in M_K$ n'est pas
triviale, ce que la page sous-entend.}. Lorsque $X$ est géométriquement
intègre, ce noyau est le faisceau constant $U_X$ défini par le groupe $U$
des $M_K$-unités de $K$ (les $u \in K^*$ tels que $v(u) = 0$ pour toute
$v \in M_K$)\footnote{Pour un corps de nombres muni de toutes ses places, $U$
est le groupe des racines de l'unité ; pour un corps de fonctions muni de
toutes ses places, le groupe multiplicatif du corps des constantes.
Remarque nôtre.}. Le même groupe est le noyau de $\mathcal{O}_X^* \to
\Phi_X$, de sorte que $\mathrm{Im}\,\mathcal{O}_X^* = \mathcal{O}_X^*/U_X$.
On a donc les suites exactes
\[
  0 \to \Phi_X \to QF_X \to \mathrm{Div}_X \to 0, \qquad
  0 \to \mathcal{R}^*_X/U_X \to QF_X \to \Phi_X/(\mathcal{O}_X^*/U_X) \to 0,
\]
qui s'assemblent, avec $0 \to \mathcal{O}_X^*/U_X \to \mathcal{R}^*_X/U_X \to
\mathrm{Div}_X \to 0$, en un diagramme commutatif à lignes et colonnes
exactes : $\Phi_X/(\mathcal{O}_X^*/U_X)$ est à la fois le conoyau de
$\mathcal{O}^*_X/U_X \to \Phi_X$ et celui de $\mathcal{R}^*_X/U_X \to QF_X$,
et $\mathrm{Div}_X$ le conoyau des deux lignes\footnote{La page dessine ce
diagramme $3 \times 3$, en partie en pointillés ; on en donne le contenu en
prose.}.

\pagerange{44}{48}
\subsection*{Sections globales, et la surjectivité sur $\mathrm{Pic}$ (pages 44, 46 et 48)}

Soit $X$ géométriquement intègre. Alors $\mathcal{R}^*_X/U_X$ et $U_X$ sont
des faisceaux constants sur un espace irréductible, donc flasques ; en
particulier $H^1(X, \mathcal{O}_X^*) \simeq H^1(X, \mathcal{O}_X^*/U_X) =
\mathrm{Pic}(X)$. En passant aux sections globales, on obtient un diagramme
commutatif à lignes et colonnes exactes, dont les trois lignes sont
\[
  \begin{array}{l}
  0 \to \Phi(X)/(\mathbf{G}_m(X)/U) \to \Gamma\bigl(X, \Phi_X/(\mathcal{O}_X^*/U_X)\bigr)
    \to \mathrm{Pic}(X) \to H^1(X, \Phi_X), \\[2pt]
  0 \to \Phi(X) \to QF(X) \to \mathrm{Div}(X) \to H^1(X, \Phi_X), \\[2pt]
  0 \to \mathbf{G}_m(X)/U \to \mathcal{R}^*(X)/U \to \mathrm{Div}_\ell(X) \to 0,
  \end{array}
\]
où $\mathrm{Div}_\ell(X)$ est le groupe des diviseurs principaux\footnote{La
phrase de la page 44 se poursuit en tête de la page 46 ; la page 45, barrée,
s'intercale dans l'ordre des archivistes (voir plus bas). L'indice de
$\mathrm{Div}_\ell$ est une lecture incertaine.}.

Il affirme : \emph{si $X$ est quasi-projectif}, les flèches vers
$H^1(X, \Phi_X)$ sont nulles ; autrement dit
\[
  \Gamma\bigl(X, \Phi_X/(\mathcal{O}_X^*/U_X)\bigr) \longrightarrow
  \mathrm{Pic}(X)
\]
est surjectif, et tout diviseur de Cartier est le diviseur d'une
quasi-fonction globale\footnote{La page dit qu'on peut « remplacer
$H^1(X, \Phi_X)$ par $0$ (deux fois) » dans le diagramme, puis
« en d'autres termes » énonce la surjectivité. Ce qui est équivalent à
celle-ci est la nullité des flèches vers $H^1$, non celle de $H^1$
lui-même ; on garde la première lecture.}. C'est la forme faisceautique de
l'existence des hauteurs locales : toute classe de faisceau inversible
porte une fonction de Weil, unique à une fonction admissible près. Tout
$\xi \in \mathrm{Pic}(X)$ étant différence de deux classes très amples, il
suffit de traiter une classe très ample ; par fonctorialité, on se ramène à
$X = \mathbf{P}^r_K$ et $\xi = \mathcal{O}(1)$ (page 48). Le feuillet
s'arrête là\footnote{Le cas de $\mathbf{P}^r$, que la page ne fait pas, se
traite ainsi (argument nôtre). Sur $U_i = \{x_i \neq 0\}$, posons
$\varphi_i(P, v) = \min_j v\bigl(x_j(P)/x_i(P)\bigr)$. Sur $U_i \cap U_j$,
$\varphi_i - \varphi_j = v\bigl(x_j/x_i\bigr)(P)$, qui est $\varphi_g$ pour
le cocycle $g = x_j/x_i$ de $\mathcal{O}(1)$ ; et $\varphi_i$ est
admissible sur $U_i$, puisque $c(v) \leqslant \varphi_i \leqslant 0$ sur
$E(U_i; (x_j/x_i)_j; c)$. Les $\varphi_i$ forment donc une section de
$\Phi/(\mathcal{O}^*/U)$ d'image $\mathcal{O}(1)$. C'est la fonction de Weil
standard, $\lambda(P, v) = \min_j v(x_j(P)) - v(x_i(P))$. La réduction par
fonctorialité demande que l'image réciproque d'une fonction admissible par
un morphisme le soit, ce que donne la remarque sur les images de parties
bornées, page 37.}.

\pagerange{45}{47}
\subsection*{Un fil barré : l'admissibilité sur les points rationnels (pages 45 et 47)}

Les pages 45 et 47 sont barrées de longs traits et se suivent entre elles.
Elles reprennent la définition sur les seuls points rationnels : une
fonction $\varphi : X(K) \times M_K \to \mathbb{R}$ est dite localement
$M_K$-bornée, ou admissible, si elle est $M_K$-bornée sur les parties
$M_K$-bornées. Alors, pour $P \in X(K)$, $\varphi(P, \cdot)$ est un
$M_K$-diviseur, et $\varphi$ s'interprète comme une application
$\tilde\varphi : X(K) \to \mathrm{Div}(M_K)$ qui envoie toute partie
$F$ telle que $F \times M_K$ soit $M_K$-bornée sur une partie bornée de
$\mathrm{Div}(M_K)$. La page demande : « Cette condition est-elle
suffisante ?? »\footnote{La question reste sans réponse sur la page, et nous
ne la tranchons pas. On remarque seulement que la condition ne porte que
sur les parties « produits » $F \times M_K$, alors qu'une partie
$M_K$-bornée n'a en général pas cette forme : pour un $X$ affine non
propre, ses fibres au-dessus des places varient.}. La page 47 ajoute que
$\tilde\varphi(X(K))$ est alors bornée, et commence le passage à une
extension finie $K'$ : une fonction sur $X(K') \times M_{K'}$ définit une
fonction sur $X(K) \times M_K$. Elle s'arrête là\footnote{Le haut de la page
45 est très surchargé et plusieurs mots en sont illisibles ; on n'en donne
que la ligne d'argument qui se lit. La page 49, dernière du dossier, est une
page dactylographiée d'un autre texte, barrée, sans rien de sa main.}.

\end{document}
